%%%%%%%%%%%%%%%%%%%%%%%%%%%%%%%%%%%%%%%%%%%%%%%%%%%%%%%%%%%%%%%%%%%%%%
% Lớp 11 | Hàm số lượng giác và phương trình lượng giác | [Giảm] Phương trình lượng giác thường gặp
% Dạng 1D1?6-5: Phương trình lượng giác không mẫu mực
%%%%%%%%%%%%%%%%%%%%%%%%%%%%%%%%%%%%%%%%%%%%%%%%%%%%%%%%%%%%%%%%%%%%%%


% nguồn: 2026-11-15_11_1D1_TL2_Luong_giac (nhập 15/11/2026)
\begin{bt}%[1D1C6-5]%[ma:751f8f9f]
% nguon: Archive/SP_KhoLatex/ID6_old/1K1-4-8.tex #3016
[VDC]%[DCHT Toán 11 - KNTT -Tên Nguyễn Thanh Sang]  
$8 \cos ^{3}\left(x+\dfrac{\pi}{3}\right)=\cos 3 x$.
\dapso{$\dfrac{\pi}{6}  + k\pi, k\pi,-\dfrac{2\pi}{3} + k\pi$}
\loigiai{
	Đặt $t=x+\dfrac{\pi}{3}$ thì $\cos 3x = \cos (3t - \pi) = -\cos 3t = 3\cos t - 4 \cot^3 t$.\\
	Khi đó phương trình trở thành:
	\begin{eqnarray*}
		8\cos^3 t = 3\cos t - 4\cos^3 t &\Leftrightarrow &3\cos t(4\cos^2 t - 1) = 0\\
		&\Leftrightarrow &\cos t(2\cos 2t + 1) = 0\\
		&\Leftrightarrow &\hoac{
			&\cos t = 0\\ &\cos 2t = -\dfrac{1}{2} = \cos\dfrac{2\pi}{3}}\\
		&\Leftrightarrow &\hoac{
			&t = \dfrac{\pi}{2} + k\pi\\ &2t = \pm\dfrac{2\pi}{3} + k2\pi
			\Leftrightarrow t = \pm\dfrac{\pi}{3} + k\pi.}
	\end{eqnarray*}
	Khi đó $x=\dfrac{\pi}{6} + k\pi$ hoặc $x=k\pi$ hoặc $x=-\dfrac{2\pi}{3} + k\pi$ với $k\in\mathbb{Z}$.\\
	Vậy tập nghiệm của phương trình là
	$S = \left\lbrace \dfrac{\pi}{6}  + k\pi, k\pi,-\dfrac{2\pi}{3} + k\pi\ \middle| \ k\in\mathbb{Z} \right\rbrace.$}
\end{bt}

% nguồn: 2026-11-15_11_1D1_TL2_Luong_giac (nhập 15/11/2026)
\begin{bt}%[1D1V6-5]%[ma:a83ca62b]
% nguon: Archive/SP_KhoLatex/ID6_old/1K1-4-10.tex #3616
%[DCHT-11-KNTT]%[Vĩ Lê Văn]
	Giải phương trình $\tan x-2\sqrt{2} \sin x=1$.
		\dapso{$\hoac{&x=\dfrac{5\pi}{12}+k 2\pi \\& x=\dfrac{13\pi}{12}+k 2\pi\\&x=-\dfrac{\pi}{4}+k 2\pi}, \quad k \in \mathbb{Z}$.}
		\loigiai{
		Điều kiện $\cos x \neq 0$.\\
		Ta có $\tan x-2\sqrt{2} \sin x=1\Leftrightarrow \sin x-\cos x-2\sqrt{2} \sin x \cos x=0$.\\
		Đặt $t=\sin x-\cos x=\sqrt{2} \sin \left(x-\dfrac{\pi}{4}\right) \Rightarrow \sin x \cos x=\dfrac{1-t^2}{2}$, điều kiện $-\sqrt{2} \leq t \leq \sqrt{2}$.\\
		PT trở thành: $t-\sqrt{2}\left(1-t^2\right)=0\Leftrightarrow \sqrt{2} t^2+t-\sqrt{2}=0\Leftrightarrow\hoac{&t=\dfrac{\sqrt{2}}{2}\quad (\text{nhận}) \\& t=-\sqrt{2}\quad (\text{nhận}).}$\\
		$\Leftrightarrow \hoac{&\sqrt{2}\sin \left(x-\dfrac{\pi}{4}\right)=\dfrac{\sqrt{2}}{2}\\&\sqrt{2}\sin \left(x-\dfrac{\pi}{4}\right)=-\sqrt{2}} \Leftrightarrow \hoac{&\sin \left(x-\dfrac{\pi}{4}\right)=\dfrac{1}{2}\\&\sin \left(x-\dfrac{\pi}{4}\right)=-1} \Leftrightarrow\hoac{&x=\dfrac{5\pi}{12}+k 2\pi \\& x=\dfrac{13\pi}{12}+k 2\pi\\&x=-\dfrac{\pi}{4}+k 2\pi}, \quad k \in \mathbb{Z}$.
		\\
		Kết hợp với điều kiện, PT có nghiệm: $\hoac{&x=\dfrac{5\pi}{12}+k 2\pi \\& x=\dfrac{13\pi}{12}+k 2\pi\\&x=-\dfrac{\pi}{4}+k 2\pi}, \quad k \in \mathbb{Z}$.}
\end{bt}
